%----------
%--- RESULTADOS

\chapter{FUNDAMENTAÇÃO TEÓRICA}
\thispagestyle{myheadings}

Este capítulo apresenta o embasamento teórico que fundamenta esta pesquisa, abordando os 
conceitos essenciais de Big Data e de integração de dados via ETL 
(\textit{Extract, Transform, Load}). Inicialmente, exploram-se a evolução do paradigma de 
Big Data, suas dimensões fundamentais (os 5 Vs) e os elementos de sua arquitetura. 
Em seguida, detalham-se as etapas e aplicações dos processos de ETL, suas variações 
(como o paradigma ELT), as principais ferramentas utilizadas pelo mercado e os aspectos 
decisivos para o projeto de pipelines eficientes de dados.

\section{Big Data}

Embora não exista uma definição única e consensual para Big Data, observa-se um certo 
consenso sobre as características que definem o Big Data. Para a Oracle, Big Data são 
dados que contêm maior variedade, chegando em volumes crescentes e com mais velocidade 
\cite{oracle_bigdata}. Já para o Google, o Big Data é um conjunto de dados maior e mais 
complexo, de modo que software tradicional não consegue processá-los e gerenciá-los 
\cite{googlecloud2025bigdata}.

Com base nessas definições, podemos observar um certo padrão para o Big Data, grande 
variedade, o alto volume e a velocidade com que os dados são gerados. Esses elementos 
convergem para uma das abordagens mais difundidas na literatura, os 3 Vs, proposto por 
Doug Laney. Nesse contexto, o Big Data é caracterizado por três dimensões fundamentais: 
volume, velocidade e variedade \cite{laney2001}.

No que se refere ao volume, este diz respeito à gigantesca quantidade dos dados gerados e 
armazenados, 90\% dos dados existentes no mundo hoje foram criados somente nos últimos dois 
anos \cite{sherman2015business}. Essa crescente no volume de dados gerados deve-se principalmente a 
fatores como a digitalização das atividades humanas, a popularização da internet e aumento 
da capacidade de armazenamento de dados.

Em relação a velocidade, ela está relacionada com a rapidez com que os dados são produzidos 
junto a necessidade de processá-los em tempo real. Atualmente dependemos da velocidade de 
alguns desses dados. "É extremamente útil receber uma notificação imediata do banco quando 
uma transação fraudulenta é detectada […]" \cite[p. 26]{sherman2015business}.

Quanto à variedade, essa característica refere-se à diversidade de fontes de origem dos dados, 
bem como aos diferentes formatos em que eles podem ser representados, como imagens, áudios, textos 
e vídeos, além de dados semiestruturados, a exemplo de documentos XML (\textit{eXtensible Markup Language}) 
e JSON (\textit{JavaScript Object Notation}). Nesse contexto, "atualmente, as empresas podem coletar 
dados de publicações em redes sociais sobre seus produtos […] assim como lidar com dados 
valiosos gerados em planilhas e documentos de texto" \cite[p. 27]{sherman2015business}.

Ademais, a definição de Big Data em termos dos 3 Vs evoluiu, ganhando dois Vs adicionais 
de veracidade e valor, formando os 5 Vs do Big Data. A veracidade, se refere à confiabilidade
e à qualidade dos dados, considerando aspectos como inconsistências, ruídos e possíveis 
incertezas presentes nas informações. Já o valor diz respeito à capacidade de extrair 
insights úteis e relevantes a partir dos dados, transformando-os em informação estratégica 
que possa apoiar a tomada de decisão e gerar benefícios para as organizações.

No mais, as características intrínsecas representadas pelos Vs do Big Data permitem 
diferenciá-lo dos modelos tradicionais de dados. Estes, por sua vez, são geralmente 
limitados a volumes menores, como gigabytes, o que torna seu armazenamento e processamento 
mais viáveis \cite{siddiqui2024bigdata}. Além disso, as fontes de dados tradicionais tendem a ser 
mais estruturadas e centralizadas, incluindo bases operacionais, registros de transações e 
sistemas de planejamento de recursos empresariais \cite{siddiqui2024bigdata}.

Diante das limitações dos modelos tradicionais de dados, incapazes de lidar com as demandas 
do Big Data, novas ferramentas e paradigmas surgiram a fim de resolver esses problemas. 
Modelos de programação como o MapReduce permitiram a organizações como a Apache 
desenvolverem todo um ecossistema voltado para solução de problemas envolvendo o Big Data, 
hoje ferramentas como o Apache Hadoop são amplamente utilizadas em empresas convencionais 
\cite{white2012hadoop}.

Desse modo, o Big Data configura-se como um paradigma caracterizado pelos 5 Vs, superando as limitações 
dos modelos tradicionais de processamento de dados. Nesse contexto, a adoção de tecnologias, como aquelas 
pertencentes ao ecossistema Apache, evidencia a necessidade de novas abordagens para o tratamento de grandes 
volumes de dados. Assim, torna-se relevante compreender o contexto de evolução dessas tecnologias e suas 
aplicações em ambientes de Big Data.

\subsection{Contextualização do Big Data}

O cenário de Big Data surgiu naturalmente como resposta ao acelerado processo de digitalização
das atividades humanas que se intensificou no início dos anos 2000. Embora, segundo 
\citeonline{gandomi2015beyond}, o conceito de Big Data provavelmente teve origem em conversas informais 
durante almoços na Silicon Graphics Inc. (SGI) em meados da década de 1990.

A partir dessas primeiras décadas, o Big Data passou a evoluir à medida que o volume, a 
velocidade e a variedade dos dados se ampliaram. Com esse crescimento, as ferramentas 
tradicionais tornaram-se insuficientes para lidar com esse novo paradigma, o que impulsionou 
o surgimento de novas soluções tecnológicas. Nesse contexto, empresas como Google e 
Microsoft, bem como a Fundação Apache, destacaram-se como importantes provedoras de 
plataformas e soluções para o cenário de Big Data \cite{khamaru2022dynamics}.

Essas novas soluções, se mostraram essenciais para o novo cenário global. Instituições que 
já dependiam de dados viram-se cada vez mais inundadas pelos grandes contingentes de 
informações. Em contrapartida, muitas empresas passaram a enxergar no cenário de Big Data 
uma oportunidade de extrair insights relevantes dos dados, com o objetivo de apoiar a tomada 
de decisões estratégicas, e assim, obter vantagens competitivas.

Nesse sentido, o Big Data passou a atuar de forma integrada às práticas de Business 
Intelligence (BI), ampliando a capacidade analítica e o uso estratégico das informações 
nas organizações. Sob essa ótica, \citeonline{narulita2025bigdata} concluem que o Big 
Data e o Business Intelligence constituem pilares estratégicos na contabilidade gerencial, 
permitindo às organizações obter vantagem competitiva por meio de decisões mais rápidas e 
precisas baseadas em dados.

\subsection{Arquitetura do Big Data}

Uma arquitetura de Big Data gerencia a ingestão, o processamento e a análise de dados muito 
grandes ou complexos para sistemas de banco de dados tradicionais \cite{microsoft2025bigdata}. 
Nesse contexto, a arquitetura organiza-se por meio de diferentes elementos e processos que, 
de forma integrada, viabilizam o tratamento eficiente dos dados ao longo de seu ciclo de vida.

\subsection{Elementos do Big Data}

A infraestrutura de armazenamento é um dos principais pilares do Big Data, é nela que o 
volume massivo e diverso dos dados é armazenado. Dado as características do Big Data, os 
principais mecanismos de armazenamento adotados são os sistemas de arquivos distribuídos, 
NoSQL, NewSQL e outros sistemas de gerenciamento de dados \cite{wang2020architecture}. 

No mais, não basta apenas armazenar dados, sendo necessário processar grandes volumes a 
fim de extrair informações relevantes. Segundo \citeonline{albarznji2022big}, o processamento de Big Data 
é essencial para transformar dados em informação útil, classificando os frameworks em cinco 
categorias principais: \textit{batch processing, streaming processing, real-time processing, 
interactive processing e hybrid processing.}

Nesse contexto, temos que o \textit{batch processing} processa grandes volumes de dados acumulados, 
enquanto o \textit{streaming processing} atua sobre fluxos contínuos de dados. 
Ja o \textit{real-time processing} prioriza respostas com baixa latência, ao passo que o 
\textit{interactive processing} permite consultas \textit{ad-hoc} (consultas formuladas sob demanda) 
e exploração dinâmica dos dados. Por fim, o \textit{hybrid processing} combina duas ou mais dessas 
abordagens para atender a diferentes demandas operacionais e analíticas.

Ademais, os dados devem ser coletados a partir de uma ou mais fontes. Nesse contexto, os 
sistemas de ingestão de dados são responsáveis pela coleta e integração dessas informações. 
Assim, ferramentas de ETL são amplamente utilizadas nesse processo, permitindo a extração, 
transformação e carregamento de diferentes tipos de dados \cite{wang2020architecture}.

Por fim, as ferramentas de Business Intelligence permitem que os usuários interajam com 
a infraestrutura de dados por meio de recursos como visualização, consultas, dashboards 
e análises preditivas \cite{sherman2015business}. Essas ferramentas são essenciais para a extração 
de insights relevantes a partir dos dados. Elas são o destino final no que tange a 
utilidade dos dados.

\subsection{Processos do Big Data}

Os processos da arquitetura de dados abrangem a preparação dos dados, o data 
franchising, as atividades de Business Intelligence (BI), bem como o gerenciamento 
de dados e de metadados. A preparação envolve atividades como coleta, transformação, 
limpeza e armazenamento. Já o data franchising refere-se à reorganização dos dados para fins analíticos. 
Por sua vez, as soluções de BI possibilitam o acesso e a análise desses dados por meio de 
ferramentas específicas. Além disso, o gerenciamento de dados e de metadados estabelece 
padrões, governança e controle sobre os ativos informacionais da organização 
\cite{sherman2015business}.

\section{ETL}

No âmbito dos processos de integração de dados, o ETL (Extract, Transform, Load) 
desempenha papel fundamental no cenário de Big Data. Segundo \citeonline{singh2022etl}, "os 
processos de ETL podem consumir até 80\% do esforço em projetos de BI [...] envolvendo a 
extração de dados de fontes externas, sua transformação para atender às necessidades de 
negócios e, por fim, seu armazenamento em um repositório de informações".

Nesse contexto, a sigla ETL refere-se às etapas de extração (Extract), transformação 
(Transform) e carregamento (Load), que, se executadas de forma estruturada, garantem a 
qualidade, consistência e integridade dos dados. Assim, para uma melhor compreensão de 
sua aplicação, torna-se necessário detalhar cada uma dessas etapas, evidenciando suas 
características, técnicas e desafios no processo de integração de dados.

A primeira etapa do processo de ETL é denominada Extract. Seu principal objetivo consiste
em capturar os dados de interesse a partir de diferentes fontes de dados 
\cite{assisvilela2021realtime}. Nessa fase, realiza-se a identificação, seleção e 
coleta dos dados provenientes de das mais diversas fontes, como bancos de dados 
relacionais, arquivos XML e aplicações externas, Essa etapa garante  que as informações 
relevantes sejam disponibilizadas para as etapas subsequentes.

Após a etapa de extração, os dados passam por um processo de transformação antes de 
serem carregados no destino final. Essa etapa é necessária devido à heterogeneidade 
das fontes e dos formatos de origem. Na fase de Transform, os dados são transformados 
ara corresponder ao esquema de dados de destino, geralmente um data warehouse 
\cite{venkateswarlu2023cloud}.

Por fim, os dados extraídos e transformados precisam ser carregados no repositório de 
destino para que possam ser efetivamente utilizados. Essa etapa, denominada Load, 
consiste na inserção dos dados no banco de dados final \cite{singh2022etl}. Trata-se da 
última fase do processo de ETL, sendo responsável por disponibilizar as informações 
de forma estruturada para diversas atividades no cenário de Big Data e BI.

\subsection{Aplicações do ETL}

Os processos de ETL apresentam ampla aplicabilidade nos mais diversos contextos, 
sendo essenciais para a integração, organização e disponibilização de informações em 
diferentes cenários. Dentre as principais aplicações do ETL, destacam-se a integração 
de dados provenientes de múltiplas fontes, processos de limpeza e qualidade de dados, 
a migração de dados e a replicação de dados.

No mais, \citeonline{arsyad2021dynamic} define a integração de dados de múltiplas fontes no contexto de 
ETL como a combinação de técnicas de engenharia e de BI usadas para combinar dados de 
diferentes fontes em informações significativas e valiosas.

Ademais, outra atividade na qual processos de ETL são comumente empregados é no data 
warehousing. O data warehouse é um banco de dados em que dados de várias origens são 
combinados para que possam ser analisados coletivamente. O ETL é geralmente usado para 
mover os dados para um data warehouse \cite{googlecloudetl}.

Outra atividade que se beneficia das etapas de um ETL, principalmente no que diz respeito
a etapa de Transform, é a limpeza e qualidade de dados. A limpeza de dados remove erros 
e mapeia os dados de origem para o formato de dados de destino
\cite{aws2026etl}.

Ademais, as ferramentas de ETL são amplamente utilizadas para migração de dados e 
para a replicação de bancos de dados. Segundo \citeonline{googlecloudetl}, as 
empresas estão movendo seus dados e aplicativos locais para a nuvem, a fim de 
economizar dinheiro, nesse contexto, o ETL é muito usado para executar essas 
migrações. No mais, a \citeonline{ibm2026etl} reforça que, o processo de data replication 
(Replicação de Dados) é frequentemente listada como um método de integração de dados que envolve ETL.

\subsection{Variações ETL e ELT}

O modelo de ETL é o mais tradicional e mais utilizado no Big Data; entretanto, não é o único. 
Existem outras variações do processo, como o ELT (\textit{Extract, Load, Transform}). 
Embora ETL e ELT sejam métodos de integração de dados, a diferença entre eles está no momento da 
transformação de dados. O ETL processa os dados transformando-os antes de carregá-los no sistema de 
destino. No ELT, os dados são carregados no sistema de destino no formato bruto e depois transformados 
\cite{googlecloudetl}.

Ademais, o modelo ELT mostra-se especialmente vantajoso em ambientes de destino com alta 
capacidade de processamento ou quando há necessidade de maior flexibilidade analítica. 
Nesse contexto, \citeonline{naphade2025evolution} destaca que o paradigma ELT aproveita o processamento massivamente
paralelo das plataformas modernas, transferindo a lógica de transformação para o ambiente de 
destino. Tal abordagem é particularmente relevante em cenários cujos requisitos analíticos 
estão em constante evolução.

\subsection{Ferramentas de ETL}

Dado o amplo uso do ETL e de seus derivados, é natural que empresas e instituições tenham 
desenvolvido ferramentas e frameworks voltados para esse tipo de processo. Nesse contexto, 
destacam-se soluções como o Apache Airflow, o SQL Server Integration Services (SSIS), o 
Apache NiFi e o Pentaho Data Integration, amplamente utilizadas na construção e gerenciamento 
de pipelines de dados.

\begin{itemize}
    \item \textbf{Apache Airflow:} é uma ferramenta de código aberto voltada à automação de fluxos de 
    trabalho. o Airflow surgiu como alternativa às ferramentas tradicionais de ETL e sendo 
    projetado para orquestrar processos computacionais complexos de forma dinâmica e escalável 
    \cite{eeti2022efficient}. Além disso, sua flexibilidade permite a integração com ferramentas de 
    processamento distribuído, como o Apache Spark. No mais, a plataforma baseia-se no conceito 
    de grafos acíclicos direcionados (DAGs) e utiliza a linguagem Python para a programação dos 
    fluxos. 

    \item \textbf{SSIS:} o SQL Server Integration Services (SSIS) é uma plataforma da Microsoft 
    voltada à integração e transformação de dados, amplamente utilizada na construção de processos de ETL. 
    Segundo \citeonline{microsoftssis}, o SSIS é utilizado para resolver problemas de negócios 
    complexos. Ademais, o SSIS se destaca como uma plataforma poderosa para executar processos 
    de ETL, permitindo que as empresas manipulem grandes volumes de dados, garantindo precisão e 
    consistência \cite{banoth2024ssis}.

    \item \textbf{Apache NIFI:} é uma ferramenta ETL de código aberto usada para automatizar fluxos que 
    lidam com dados em tempo real. O NiFi foi desenvolvido pela Apache Software Foundation e é 
    adequado para tarefas de fluxo de dados, processamento e interação com sistemas, tudo isso 
    apresentado por meio de um sofisticado fluxograma gráfico \cite{singu2022automation}. 

    \item \textbf{Pentaho:} é uma plataforma de Business Intelligence amplamente utilizada para 
    integração, transformação e análise de dados. A ferramenta possibilita a execução de 
    processos de ETL, além de oferecer suporte à criação de relatórios, dashboards e outros 
    mecanismos de visualização. Segundo \citeonline{demussa2018business}, o Pentaho dispõe de uma 
    infraestrutura completa para a realização de operações de BI, abrangendo desde ferramentas 
    de ETL até mecanismos de acesso, controle e suporte à elaboração de relatórios e painéis 
    analíticos.
\end{itemize}

\subsection{Projeto de ETL}

O projeto de um ETL consiste na definição de estratégias e técnicas para a coleta, 
tratamento e disponibilização de dados provenientes de múltiplas fontes. Essa etapa é 
fundamental, pois garante que os dados sejam integrados, consistentes e adequados para 
análises. Um projeto de ETL bem estruturado deve considerar aspectos como qualidade dos dados, 
desempenho, escalabilidade e governança, assegurando que as informações estejam disponíveis de 
forma confiável para apoio à tomada de decisão.

\begin{itemize}
    \item \textbf{Data Extraction:} a etapa de extração de dados (Data Extraction) é responsável por 
    coletar dados de diferentes fontes, como bancos de dados relacionais, arquivos, APIs e 
    sistemas, planilhas etc. Essa fase é crítica, pois impacta diretamente a qualidade e a 
    integridade das etapas subsequentes do processo de ETL. A extração pode ocorrer de forma 
    completa ou incremental, sendo esta última utilizada para capturar apenas dados novos ou 
    modificados, reduzindo o volume de processamento e aumentando a eficiência do sistema. 
    Além disso, é importante que o processo de extração minimize impactos nos sistemas de 
    origem, garantindo a continuidade das operações.

    \item \textbf{Data Franchising:} o data franchising refere-se ao processo de disponibilização dos 
    dados já tratados em repositórios voltados ao consumo pelos usuários. Essa abordagem 
    consiste em empacotar os dados em estruturas como data marts (subconjuntos de dados 
    orientados a áreas específicas) ou cubos analíticos, facilitando o acesso, a compreensão 
    e a utilização das informações por analistas e gestores. Embora essa prática gera 
    redundância em relação ao data warehouse, trata-se de uma redundância controlada e 
    intencional, que visa melhorar o desempenho e a usabilidade dos dados. Ademais, o data 
    franchising ocorre após as etapas de preparação e transformação, garantindo que apenas 
    dados consistentes e estruturados sejam disponibilizados \cite{sherman2015business}.

\end{itemize}














